% ============================================
% FILE: run_length_encoding.tex
% Detailed Lecture Notes for Chapter 7: Run-Length Encoding
% ============================================

\section{Run-Length Encoding: Exploiting Equality Redundancy}

\begin{importantbox}
\textbf{Learning Objectives:}
\begin{itemize}
    \item Understand the concept of equality redundancy and why it's the simplest form of compressible pattern
    \item Master the basic RLE algorithm and its variants for different data types
    \item Analyze compression ratios, best/worst cases, and theoretical limits of RLE
    \item Learn four major RLE variants: Count-Based, Escape-Based, Block-Based, and Zero RLE
    \item Understand why RLE fails on some data types and how to combine it with other techniques
    \item Build the conceptual bridge to transform-based compression, especially BWT
\end{itemize}
\end{importantbox}

\subsection{Introduction and Motivation}

\subsubsection{The Simplest Form of Redundancy: Repetition}

When we look at data through the lens of compression, we're constantly asking: "What patterns repeat?" Among all possible patterns, the absolute simplest is when the \textit{same symbol repeats consecutively}. We call this \textbf{equality redundancy}.

\begin{definitionbox}
\textbf{Equality Redundancy:} A form of data redundancy where the same symbol appears multiple times in consecutive positions within a sequence. Each such group of identical consecutive symbols is called a \textbf{run}.
\end{definitionbox}

Consider these examples:
\begin{itemize}
    \item Text: \texttt{AAAABBBCCCC} contains three runs: \texttt{AAAA}, \texttt{BBB}, and \texttt{CCCC}
    \item Images: A blue sky photo has long runs of similar blue pixel values
    \item Signals: Silence in audio contains runs of zero amplitude
    \item Documents: White space in scanned text creates runs of white pixels
\end{itemize}

The key insight is elegant in its simplicity: instead of storing each symbol individually, we can store each symbol once along with a count of how many times it repeats. For \texttt{AAAA}, instead of storing 4 bytes, we store (\texttt{A},4).

\begin{examplebox}
\textbf{Intuition Building:}
Raw data: \texttt{WWWWWWWWWWWWWBBWWWWWWWWWWWWWBBBWWWWWWWWW}
\newline
With runs identified: 
\begin{itemize}
    \item 13 W's
    \item 2 B's
    \item 13 W's
    \item 3 B's
    \item 9 W's
\end{itemize}
Storing as pairs: (W,13)(B,2)(W,13)(B,3)(W,9)

Original: 40 characters
Encoded: 10 values (5 symbols + 5 counts) $\rightarrow$ significant savings!
\end{examplebox}

\subsection{Fundamental Concepts and Theory}

\subsubsection{Formal Definition of a Run}

Before we can encode runs, we need a precise mathematical definition.

\begin{definitionbox}
\textbf{Run:} Given a sequence $S = s_1, s_2, s_3, \ldots, s_n$ over an alphabet $\Sigma$, a run is a maximal consecutive subsequence $s_i, s_{i+1}, \ldots, s_j$ such that $s_i = s_{i+1} = \ldots = s_j$, and either $i=1$ or $s_{i-1} \neq s_i$, and either $j=n$ or $s_{j+1} \neq s_j$.

In simpler terms: a run is the longest possible sequence of identical symbols that appears consecutively.
\end{definitionbox}

\textbf{Properties of Runs:}
\begin{enumerate}
    \item \textbf{Maximality:} A run cannot be extended—the symbol before it (if any) is different, and the symbol after it (if any) is different.
    \item \textbf{Disjointness:} Runs partition the sequence; every symbol belongs to exactly one run.
    \item \textbf{Alternation:} Adjacent runs always contain different symbols.
\end{enumerate}

\subsubsection{Run-Length Distribution Theorem}

Understanding the statistical properties of runs is fundamental to compression analysis.

\begin{theorem}
For an i.i.d. (independent and identically distributed) random sequence over an alphabet of size $k$, the expected run length is:
\[
E[\text{run length}] = \frac{1}{1-p} = \frac{k}{k-1}
\]
where $p = 1/k$ is the probability of a symbol repeating.
\end{theorem}

\begin{proof}
For a random sequence with equiprobable symbols:
\begin{itemize}
    \item Probability that next symbol equals current = $1/k$
    \item Run length follows geometric distribution: $P(\text{run length} = L) = (1/k)^{L-1} \cdot (1 - 1/k)$
    \item Expected value: $\sum_{L=1}^{\infty} L \cdot (1/k)^{L-1} \cdot (1 - 1/k) = \frac{1}{1-1/k} = \frac{k}{k-1}$
\end{itemize}
\end{proof}

\begin{examplebox}
\textbf{Expected Run Lengths for Different Alphabets:}

\begin{tabular}{|l|c|c|}
\hline
\textbf{Alphabet} & \textbf{Size ($k$)} & \textbf{Expected run length} \\
\hline
Binary & 2 & $2/(2-1) = 2$ \\
DNA bases & 4 & $4/3 \approx 1.333$ \\
Lowercase letters & 26 & $26/25 = 1.04$ \\
ASCII & 256 & $256/255 \approx 1.0039$ \\
\hline
\end{tabular}

\textbf{Implications:}
\begin{itemize}
    \item Binary data has the longest expected runs (2)
    \item ASCII text has runs barely longer than 1
    \item For large alphabets, runs of length $>$ 1 are rare events
\end{itemize}
\end{examplebox}

\subsubsection{Compression Ratio Analysis}

\begin{definitionbox}
\textbf{Compression Ratio:} 
\[
\text{Ratio} = \frac{\text{Original Size}}{\text{Compressed Size}}
\]
Higher ratio means better compression.
\end{definitionbox}

\begin{theorem}
For a sequence of length $n$ consisting of a single symbol, the best-case compression ratio achievable with basic RLE is approximately $\frac{n}{c}$, where $c$ is the size (in bytes) of one (symbol, count) pair.
\end{theorem}

\begin{theorem}
For a sequence of length $n$ with no repeated consecutive symbols (all runs of length 1), basic RLE achieves a compression ratio of $\frac{n}{2n} = 0.5$ when using fixed-size counts equal to symbol size, resulting in 100\% expansion.
\end{theorem}

\begin{proof}
If every symbol is different from its predecessor:
\begin{itemize}
    \item Number of runs = $n$ (each symbol is its own run)
    \item Encoded output = $n$ (symbol, count) pairs
    \item If count field size equals symbol size (e.g., 1 byte each):
    \begin{itemize}
        \item Original size = $n$ bytes
        \item Compressed size = $n$ symbols + $n$ counts = $2n$ bytes
        \item Ratio = $n/(2n) = 0.5$
        \item Expansion = $(2n - n)/n = 1 = 100\%$
    \end{itemize}
\end{itemize}
\end{proof}

\subsection{Taxonomy of RLE Methods}

RLE variants differ in their answers to three fundamental questions:
\begin{enumerate}
    \item How do we distinguish between runs and non-repeated data?
    \item How do we represent the count?
    \item How do we handle maximum run lengths?
\end{enumerate}

Based on these choices, we can identify four major RLE families:

\begin{importantbox}
\textbf{The Four Major RLE Variants:}

\begin{enumerate}
    \item \textbf{Count-Based RLE (Classic RLE):} Stores each run as (symbol, count). Simple and effective for data with many runs.
    
    \item \textbf{Escape-Based RLE:} Uses special markers to distinguish runs from literals. Handles mixed data efficiently.
    
    \item \textbf{Block/Header-Based RLE (PackBits Style):} Divides data into chunks with headers that describe literal or run data. Balances overhead and efficiency.
    
    \item \textbf{Zero Run-Length Encoding (Zero RLE):} Specialized for compressing runs of zeros in sparse data, commonly used in transform coding.
\end{enumerate}
\end{importantbox}

\newpage
\subsection{Count-Based RLE (Classic RLE)}

Count-based RLE is the original and simplest form—each run is represented by its symbol and a count of how many times it repeats.

\subsubsection{Basic Encoding Algorithm}

\begin{algorithm}[H]
\caption{Count-Based RLE Encoding}
\begin{algorithmic}[1]
\Procedure{RLE-Encode}{$S[1..n]$}
    \State $result \gets$ empty list
    \If {$n = 0$}
        \State \Return $result$
    \EndIf
    \State $current\_symbol \gets S[1]$
    \State $run\_length \gets 1$
    \For {$i \gets 2$ to $n$}
        \If {$S[i] = current\_symbol$}
            \State $run\_length \gets run\_length + 1$
        \Else
            \State append $(current\_symbol, run\_length)$ to $result$
            \State $current\_symbol \gets S[i]$
            \State $run\_length \gets 1$
        \EndIf
    \EndFor
    \State append $(current\_symbol, run\_length)$ to $result$ \Comment{Don't forget the last run!}
    \State \Return $result$
\EndProcedure
\end{algorithmic}
\end{algorithm}

\subsubsection{Basic Decoding Algorithm}

\begin{algorithm}[H]
\caption{Count-Based RLE Decoding}
\begin{algorithmic}[1]
\Procedure{RLE-Decode}{$pairs[1..m]$}
    \State $result \gets$ empty list
    \For {each $(symbol, count)$ in $pairs$}
        \For {$i \gets 1$ to $count$}
            \State append $symbol$ to $result$
        \EndFor
    \EndFor
    \State \Return $result$
\EndProcedure
\end{algorithmic}
\end{algorithm}

\subsubsection{Count-After-Symbol vs. Count-Before-Symbol}

The ordering of symbol and count in the output stream affects decoding and has subtle implications.

\begin{importantbox}
\textbf{Two Common Conventions:}

\textbf{1. Count-After-Symbol (Symbol then Count):}
\begin{itemize}
    \item Format: \texttt{[symbol][count][symbol][count]...}
    \item Example: \texttt{A5B3C2} means "5 A's, 3 B's, 2 C's"
    \item \textbf{Advantages:}
    \begin{itemize}
        \item Natural reading order: "Here's what to output, then how many"
        \item Easier for streaming: see symbol first, can prepare output
        \item Common in many implementations
    \end{itemize}
\end{itemize}

\textbf{2. Count-Before-Symbol (Count then Symbol):}
\begin{itemize}
    \item Format: \texttt{[count][symbol][count][symbol]...}
    \item Example: \texttt{5A3B2C} means "5 A's, 3 B's, 2 C's"
    \item \textbf{Advantages:}
    \begin{itemize}
        \item Decoder knows length before seeing symbol
        \item Can allocate buffer space if needed
        \item Used in some formats (e.g., PCX)
    \end{itemize}
\end{itemize}
\end{importantbox}

\subsubsection{Fixed-Width Count RLE}

In fixed-width count RLE, each run length is stored in a fixed number of bits.

\begin{importantbox}
\textbf{Fixed-Width Count Fields Design:}

\textbf{Design:} Each run length stored in a fixed number of bits (e.g., 8 bits, 16 bits, 32 bits)

\textbf{Advantages:}
\begin{itemize}
    \item Simple to implement
    \item Fast encoding/decoding (no bit-level parsing)
    \item Random access possible (know where each run ends)
    \item Good for hardware implementation
\end{itemize}

\textbf{Disadvantages:}
\begin{itemize}
    \item Wasted space for short runs
    \item Limited maximum run length (e.g., 255 for 8 bits)
    \item Need to split longer runs into multiple chunks
    \item Not adaptive to run length distribution
\end{itemize}

\textbf{Typical trade-offs:}

\begin{tabular}{|l|c|c|c|}
\hline
\textbf{Count size} & \textbf{Max run} & \textbf{Overhead/run} & \textbf{Best for} \\
\hline
4 bits (nibble) & 15 & 0.5 bytes & Very short runs \\
8 bits (byte) & 255 & 1 byte & General purpose \\
16 bits & 65,535 & 2 bytes & Long runs \\
24 bits & 16.7M & 3 bytes & Images \\
32 bits & 4.3B & 4 bytes & Very large files \\
\hline
\end{tabular}
\end{importantbox}

\begin{examplebox}
\textbf{Fixed-Width RLE Example:}

Input sequence: \texttt{WWWWWWWWWWWWWBBWWWWWWWWWWWWWBBBWWWWWWWWW}

\textbf{With 8-bit counts:}
\begin{itemize}
    \item Runs: (W,13), (B,2), (W,13), (B,3), (W,9)
    \item Each run: 1 byte symbol + 1 byte count = 2 bytes
    \item Total: 5 runs × 2 bytes = 10 bytes
    \item Original: 40 bytes
    \item Compression ratio: 4:1 (75\% savings)
\end{itemize}

\textbf{With 4-bit counts (nibbles):}
\begin{itemize}
    \item Need to pack two counts per byte
    \item Runs: (W,13), (B,2), (W,13), (B,3), (W,9)
    \item Symbols: 5 bytes
    \item Counts: 5 × 4 bits = 20 bits = 2.5 bytes
    \item Total: 7.5 bytes (assuming we can pack)
    \item But 13 exceeds 4-bit max (15) — OK here
    \item Compression ratio: 5.33:1 (81\% savings)
\end{itemize}
\end{examplebox}

\subsubsection{Handling Large Runs with Fixed-Width Counts}

When runs exceed the maximum representable count, we need splitting strategies.

\begin{importantbox}
\textbf{Splitting Runs Strategy:}

When count exceeds maximum, output multiple run chunks:
\begin{itemize}
    \item Example: max count = 255, actual run = 1000
    \item Output: (sym,255), (sym,255), (sym,255), (sym,235)
    \item Overhead increases but preserves maximum limit
\end{itemize}

\textbf{Alternative: Use larger count fields}
\begin{itemize}
    \item Switch to 16-bit or 32-bit counts
    \item But overhead increases for all runs
    \item May still need splitting for very large runs
\end{itemize}
\end{importantbox}

\begin{examplebox}
\textbf{Run Splitting Example:}

Assume 8-bit counts (max = 255)

Input run: 1000 'A's

\textbf{Split into max-sized chunks:}
\begin{itemize}
    \item (A,255), (A,255), (A,255), (A,235)
    \item Overhead: 4 pairs instead of 1
    \item Extra overhead: 3 extra count bytes + 3 extra symbol bytes = 6 extra bytes
    \item Still much better than storing 1000 raw bytes!
\end{itemize}
\end{examplebox}

\newpage
\subsubsection{Variable-Width Count RLE (Elias Gamma)}

Variable-width count RLE uses codes that use fewer bits for small numbers and more bits for large numbers, matching the information content of the count.

\begin{importantbox}
\textbf{Variable-Length Count Fields Design:}

\textbf{Design:} Use a variable-length code (like Elias gamma coding, Fibonacci coding) to represent run lengths

\textbf{Common approaches:}
\begin{enumerate}
    \item \textbf{Gamma/Elias codes:} Good for small numbers, inefficient for large
    \item \textbf{Fibonacci codes:} Better balance, no need for separator
    \item \textbf{Exponential-Golomb:} Used in H.264/AVC video coding
    \item \textbf{Adaptive Huffman:} Best compression, but complex
\end{enumerate}

\textbf{Advantages:}
\begin{itemize}
    \item Optimal for the run length distribution
    \item No arbitrary maximum run length
    \item No wasted bits for short runs
    \item Can achieve near-entropy compression of run lengths
\end{itemize}

\textbf{Disadvantages:}
\begin{itemize}
    \item Complex implementation (bit-level I/O)
    \item Slower encoding/decoding
    \item No random access without decoding entire stream
    \item Error propagation if bit errors occur
\end{itemize}
\end{importantbox}

\subsubsection*{Elias Gamma Coding Deep Dive}

Elias gamma coding provides the clearest illustration of variable-length coding principles.

\begin{definitionbox}
\textbf{Elias Gamma Code:} A variable-length code for positive integers where each number $N \ge 1$ is encoded as:
\begin{enumerate}
    \item Write $N$ in binary.
    \item Let the binary representation of $N$ use $L$ bits.
    \item Write $L-1$ zeros.
    \item Append the full binary representation of $N$.
\end{enumerate}

Equivalently, if $k = \lfloor \log_2 N \rfloor$, then:
\[
\text{Gamma}(N) = \underbrace{00\ldots0}_{k \text{ zeros}} \; \text{binary}(N)
\]

Total code length: $2k + 1$ bits.
\end{definitionbox}

\subsubsection*{Building Elias Gamma Codes: Step by Step}

\begin{enumerate}
    \item \textbf{For N = 1:}
    \begin{itemize}
        \item Binary: $1$ (L=1)
        \item Zeros: L-1 = 0 zeros
        \item Code: "1"
    \end{itemize}
    
    \item \textbf{For N = 2:}
    \begin{itemize}
        \item Binary: $10$ (L=2)
        \item Zeros: L-1 = 1 zero
        \item Code: "010"
    \end{itemize}
    
    \item \textbf{For N = 3:}
    \begin{itemize}
        \item Binary: $11$ (L=2)
        \item Zeros: 1 zero
        \item Code: "011"
    \end{itemize}
    
    \item \textbf{For N = 4:}
    \begin{itemize}
        \item Binary: $100$ (L=3)
        \item Zeros: 2 zeros
        \item Code: "00100"
    \end{itemize}
    
    \item \textbf{For N = 5:}
    \begin{itemize}
        \item Binary: $101$ (L=3)
        \item Zeros: 2 zeros
        \item Code: "00101"
    \end{itemize}
\end{enumerate}

\subsubsection*{Complete Reference Table}

\begin{tabular}{|c|c|c|c|}
\hline
\textbf{N} & \textbf{Binary of N} & \textbf{k = floor(log2 N)} & \textbf{Gamma Code} \\
\hline
1  & 1    & 0 & 1 \\
2  & 10   & 1 & 010 \\
3  & 11   & 1 & 011 \\
4  & 100  & 2 & 00100 \\
5  & 101  & 2 & 00101 \\
6  & 110  & 2 & 00110 \\
7  & 111  & 2 & 00111 \\
8  & 1000 & 3 & 0001000 \\
9  & 1001 & 3 & 0001001 \\
10 & 1010 & 3 & 0001010 \\
\hline
\end{tabular}

\subsubsection*{Decoding Elias Gamma}

\begin{algorithm}[H]
\caption{Elias Gamma Decoding}
\begin{algorithmic}[1]
\Procedure{DecodeGamma}{$bitstream$}
    \State $k \gets 0$
    \While {next bit is 0} \Comment{Count leading zeros}
        \State $k \gets k + 1$
        \State advance bitstream
    \EndWhile
    \State advance bitstream \Comment{Skip the 1}
    \State $value \gets 1$ \Comment{The leading 1 is the most significant bit}
    \For {$i \gets 1$ to $k$}
        \State $value \gets (value \ll 1) |$ next bit
    \EndFor
    \State \Return $value$
\EndProcedure
\end{algorithmic}
\end{algorithm}

\begin{examplebox}
\textbf{Decoding Example:}

Given bitstream: \texttt{001011...}

\begin{enumerate}
    \item Read zeros: "00" then a 1 → found "001", so $k = 2$
    \item The 1 we read is part of the code
    \item Read next $k=2$ bits: "01"
    \item Construct value: leading 1 + "01" = "101" binary = 5
    \item First number decoded is 5, remaining bits continue from "..."
\end{enumerate}
\end{examplebox}

\newpage
\subsubsection*{Fixed-Width vs. Variable-Width: Detailed Comparison}

\begin{examplebox}
\textbf{Count Field Comparison:}

Consider a file with runs of lengths: 3, 5, 2, 100, 4, 200

\textbf{Fixed 8-bit counts:}
\begin{itemize}
    \item Each run: 1 byte for count
    \item Total count overhead: 6 bytes = 48 bits
    \item Can represent all runs (max 200 $<$ 255)
    \item Simple, fast, randomly accessible
\end{itemize}

\textbf{Fixed 4-bit counts:}
\begin{itemize}
    \item Cannot represent runs $>$15
    \item Would need to split run 100 into: 15,15,15,15,15,15,10 (7 runs!)
    \item Run 200: even worse, would need 14 runs (13×15 + 5)
    \item Total overhead: 6 runs become 6+7+14 = 27 runs!
    \item 27 runs × 0.5 bytes = 13.5 bytes (worse than 8-bit!)
\end{itemize}

\textbf{Variable-length (Elias gamma):}
\begin{itemize}
    \item 3 → "011" (3 bits)
    \item 5 → "00101" (5 bits)
    \item 2 → "010" (3 bits)
    \item 100 → Let's construct: binary 100 = 1100100 (7 bits), k=6, so "0000001" + "100100" = "0000001100100" (13 bits)
    \item 4 → "00100" (5 bits)
    \item 200 → binary 200 = 11001000 (8 bits), k=7, so "00000001" + "1001000" = "000000011001000" (15 bits)
    \item Total: 44 bits $\approx$ 5.5 bytes
\end{itemize}

\textbf{Comparison Table:}

\begin{tabular}{|l|c|c|}
\hline
\textbf{Method} & \textbf{Total Bits} & \textbf{Characteristics} \\
\hline
Fixed 4-bit & Not possible (runs too long) & Would need run splitting \\
Fixed 8-bit & 48 bits & Simple, fast, random access \\
Elias gamma & 44 bits & 8.3\% savings, complex, no random access \\
\hline
\end{tabular}

\textbf{Bit-by-bit breakdown:}

\begin{tabular}{|l|c|c|c|c|c|c|c|}
\hline
\textbf{Method} & \textbf{3} & \textbf{5} & \textbf{2} & \textbf{100} & \textbf{4} & \textbf{200} & \textbf{Total} \\
\hline
Fixed 8-bit & 8 & 8 & 8 & 8 & 8 & 8 & 48 \\
Elias gamma & 3 & 5 & 3 & 13 & 5 & 15 & 44 \\
Savings & +5 & +3 & +5 & -5 & +3 & -7 & +4 \\
\hline
\end{tabular}

Notice the economic trade-off: small numbers (2,3,4,5) use fewer bits than fixed-width, which subsidizes the extra bits needed for large numbers (100,200). In typical data where short runs are much more common than long runs, this trade-off yields net compression.
\end{examplebox}

\subsubsection*{Decision Matrix: Fixed vs. Variable-Width Counts}

\begin{tabular}{|p{3.5cm}|p{5.5cm}|p{5.5cm}|}
\hline
\textbf{Design Question} & \textbf{Fixed-Width Counts} & \textbf{Variable-Width Counts} \\
\hline
Implementation complexity & Very simple. Each count always uses the same number of bits. & More complex. Must read bits carefully to determine where each count ends. \\
\hline
Encoding/decoding speed & Very fast. No need to examine individual bits. & Slower. Must count prefix bits and parse variable-length codes. \\
\hline
Random access & Yes. If each count uses (say) 8 bits, boundaries are known immediately. & No. Must decode sequentially to know where the next count begins. \\
\hline
Space efficiency & When run lengths are roughly similar or bounded (e.g., always between 1 and 200). & When most run lengths are small but occasionally very large. Small values use very few bits. \\
\hline
Maximum run length & Limited by chosen width. For 8 bits, maximum is 255. & No fixed upper bound. Large values simply use more bits. \\
\hline
Error resilience & Damage usually affects only that run. & Damage may affect all following runs because boundaries shift. \\
\hline
\end{tabular}

\newpage
\subsection{Escape-Based RLE}

Escape-based RLE addresses the fundamental weakness of count-based RLE: expansion on non-repetitive data. By using special markers, it treats runs and literals differently.

\subsubsection{Core Concept}

\begin{definitionbox}
\textbf{Escape Symbol:} A special marker that indicates whether the following data represents a run or a sequence of literal (uncompressed) symbols.
\end{definitionbox}

The key insight: instead of encoding every symbol as a (symbol, count) pair (which doubles the size for non-repeated symbols), we only encode runs when they're beneficial, and leave non-repeated data as-is.

\subsubsection{Basic Escape Scheme Design}

\begin{importantbox}
\textbf{Standard Escape Scheme:}

\begin{enumerate}
    \item Choose an escape symbol (often a value not likely to appear in data, or one that can be escaped itself)
    \item When encountering a run of length $\geq 3$:
    \begin{itemize}
        \item Output escape symbol
        \item Output the repeated symbol
        \item Output the run length
    \end{itemize}
    \item For non-repetitive data:
    \begin{itemize}
        \item Output symbols literally
        \item If the literal data contains the escape symbol, need a way to represent it (e.g., double escape)
    \end{itemize}
\end{enumerate}
\end{importantbox}

\subsubsection{Encoding Algorithm}

\begin{algorithm}[H]
\caption{Escape-Based RLE Encoding (Threshold T=3)}
\begin{algorithmic}[1]
\Procedure{EscapeRLE-Encode}{$S[1..n]$, $escape$}
    \State $result \gets$ empty list
    \State $i \gets 1$
    \While {$i \leq n$}
        \State $run\_length \gets 1$
        \While {$i + run\_length \leq n$ and $S[i + run\_length] = S[i]$}
            \State $run\_length \gets run\_length + 1$
        \EndWhile
        \If {$run\_length \geq 3$} \Comment{Long run - encode with escape}
            \State append $escape$ to $result$
            \State append $S[i]$ to $result$
            \State append $run\_length$ to $result$
        \Else \Comment{Short run - output as literals}
            \For {$j \gets 0$ to $run\_length-1$}
                \If {$S[i+j] = escape$}
                    \State append $escape$ to $result$ \Comment{Escape the escape!}
                    \State append $escape$ to $result$
                \Else
                    \State append $S[i+j]$ to $result$
                \EndIf
            \EndFor
        \EndIf
        \State $i \gets i + run\_length$
    \EndWhile
    \State \Return $result$
\EndProcedure
\end{algorithmic}
\end{algorithm}

\subsubsection{Decoding Algorithm}

\begin{algorithm}[H]
\caption{Escape-Based RLE Decoding}
\begin{algorithmic}[1]
\Procedure{EscapeRLE-Decode}{$data[1..m]$, $escape$}
    \State $result \gets$ empty list
    \State $i \gets 1$
    \While {$i \leq m$}
        \If {$data[i] = escape$}
            \If {$i+2 > m$} \State error \EndIf
            \State $symbol \gets data[i+1]$
            \State $count \gets data[i+2]$
            \For {$j \gets 1$ to $count$}
                \State append $symbol$ to $result$
            \EndFor
            \State $i \gets i + 3$
        \Else
            \If {$i+1 \leq m$ and $data[i+1] = escape$} \Comment{Double escape means literal escape}
                \State append $escape$ to $result$
                \State $i \gets i + 2$
            \Else
                \State append $data[i]$ to $result$
                \State $i \gets i + 1$
            \EndIf
        \EndIf
    \EndWhile
    \State \Return $result$
\EndProcedure
\end{algorithmic}
\end{algorithm}

\subsubsection{Worked Example}

\begin{examplebox}
\textbf{Escape Scheme Example with Escape = 0xFF:}

Input: \texttt{AAAAABCDEFFFFGG} (15 bytes)

\textbf{Step-by-step encoding:}
\begin{itemize}
    \item Position 1-5: \texttt{AAAAA} (5 A's) → run length 5 $\ge$ 3
    \begin{itemize}
        \item Output: \texttt{FF A 05}
    \end{itemize}
    \item Position 6: \texttt{B} (literal) → output: \texttt{B}
    \item Position 7: \texttt{C} (literal) → output: \texttt{C}
    \item Position 8: \texttt{D} (literal) → output: \texttt{D}
    \item Position 9: \texttt{E} (literal) → output: \texttt{E}
    \item Position 10-13: \texttt{FFFF} (4 F's) → run length 4 $\ge$ 3
    \begin{itemize}
        \item Output: \texttt{FF F 04}
    \end{itemize}
    \item Position 14-15: \texttt{GG} (2 G's) → run length 2 $<$ 3, output as literals
    \begin{itemize}
        \item Output: \texttt{G G}
    \end{itemize}
\end{itemize}

Encoded: \texttt{FF A 05 B C D E FF F 04 G G}

\textbf{Analysis:}
\begin{itemize}
    \item Original: 15 bytes
    \item Encoded: 11 bytes
    \item Savings: 4 bytes (27\% compression)
    \item Without escape scheme (basic RLE): would be 16 bytes (expansion!)
\end{itemize}
\end{examplebox}

\begin{examplebox}
\textbf{Handling Escape Symbols in Literal Data:}

Input: \texttt{ABC\textbackslash xFF\textbackslash xFFDEF} (contains two 0xFF bytes)

\textbf{Encoding:}
\begin{itemize}
    \item \texttt{A} → literal A
    \item \texttt{B} → literal B
    \item \texttt{C} → literal C
    \item \verb|\xFF| (first) → this is the escape symbol! Need to escape it
    \begin{itemize}
        \item Output: \texttt{FF FF} (double escape means literal FF)
    \end{itemize}
    \item \verb|\xFF| (second) → another escape symbol
    \begin{itemize}
        \item Output: \texttt{FF FF}
    \end{itemize}
    \item \texttt{D} → literal D
    \item \texttt{E} → literal E
    \item \texttt{F} → literal F
\end{itemize}

Encoded: \texttt{A B C FF FF FF FF D E F}

Notice how each literal 0xFF becomes two bytes in the output (the escape marker twice).
\end{examplebox}

\subsubsection{Choosing the Threshold}

The threshold determines the minimum run length that triggers escape encoding.

\begin{importantbox}
\textbf{Threshold Trade-offs:}

\begin{itemize}
    \item \textbf{Low threshold (T=2):}
    \begin{itemize}
        \item Catches more runs (even pairs)
        \item Overhead: each run costs 3 bytes instead of 2 bytes for literals
        \item Break-even: 3 bytes for run vs. 2 bytes for literals → not beneficial!
        \item T=2 actually causes expansion on pairs
    \end{itemize}
    
    \item \textbf{Medium threshold (T=3):}
    \begin{itemize}
        \item Run of 3: 3 bytes encoded vs. 3 bytes literals → break-even
        \item Run of 4+: 3 bytes encoded vs. 4+ bytes literals → savings
        \item Most common choice
    \end{itemize}
    
    \item \textbf{High threshold (T=4):}
    \begin{itemize}
        \item Only encodes longer runs
        \item Misses runs of 3 (which are break-even anyway)
        \item Safer but less aggressive
    \end{itemize}
\end{itemize}

\textbf{Optimal threshold analysis:}
With 1-byte escape, 1-byte symbol, 1-byte count:
\begin{itemize}
    \item Cost to encode run of length L: 3 bytes
    \item Cost to store as literals: L bytes
    \item Break-even when L = 3
    \item Savings when L $>$ 3: L - 3 bytes saved
    \item Loss when L $<$ 3: 3 - L bytes wasted
\end{itemize}
\end{importantbox}

\subsubsection{Advantages and Disadvantages}

\begin{importantbox}
\textbf{Escape-Based RLE: Pros and Cons}

\textbf{Advantages:}
\begin{itemize}
    \item No expansion on non-repetitive data (except for escaping the escape symbol)
    \item Handles mixed data types efficiently
    \item Simple to implement
    \item Used in many real-world formats (e.g., TIFF, some fax formats)
\end{itemize}

\textbf{Disadvantages:}
\begin{itemize}
    \item Overhead of escape symbol reduces efficiency on highly repetitive data
    \item Need to handle escape symbol in literal data (double-escape)
    \item Fixed threshold may not be optimal for all data
    \item 3-byte overhead per run (vs. 2 bytes in basic RLE)
\end{itemize}
\end{importantbox}

\newpage
\subsection{Block/Header-Based RLE (PackBits Style)}

Block-based RLE divides the data into chunks, each preceded by a header that describes how to interpret the following bytes. PackBits is the classic example of this approach.

\subsubsection{PackBits Overview}

\begin{definitionbox}
\textbf{PackBits Encoding:} Each chunk begins with a 1-byte header $h$ interpreted as a signed 8-bit integer:
\begin{itemize}
    \item If $0 \leq h \leq 127$: copy the next $h+1$ bytes literally
    \item If $-127 \leq h \leq -1$: repeat the next byte $1-h$ times
    \item $h = -128$: reserved (no operation, sometimes used for special purposes)
\end{itemize}
\end{definitionbox}

\begin{importantbox}
\textbf{PackBits Header Interpretation:}

\textbf{Positive (0 to 127):} Literal run
\begin{itemize}
    \item Value $n$ means "copy next $n+1$ bytes as-is"
    \item Range: 1 to 128 bytes per literal chunk
\end{itemize}

\textbf{Negative (-1 to -127):} Repeated run
\begin{itemize}
    \item Value $n$ means "repeat next byte $1-n$ times"
    \item Since $n$ is negative, $1-n$ is positive
    \item Range: 2 to 128 repetitions
    \item Example: $h = -3$ means repeat next byte $1-(-3) = 4$ times
\end{itemize}

\textbf{-128:} Reserved (some implementations ignore, others use as EOF marker)
\end{importantbox}

\subsubsection{Encoding Algorithm}

\begin{algorithm}[H]
\caption{PackBits Encoding}
\begin{algorithmic}[1]
\Procedure{PackBits-Encode}{$S[1..n]$}
    \State $result \gets$ empty list
    \State $i \gets 1$
    \While {$i \leq n$}
        \State $run\_length \gets 1$
        \While {$i + run\_length \leq n$ and $S[i + run\_length] = S[i]$ and $run\_length < 128$}
            \State $run\_length \gets run\_length + 1$
        \EndWhile
        
        \If {$run\_length \geq 2$} \Comment{We have a run}
            \State $header \gets -(run\_length - 1)$ \Comment{Negative header}
            \State append $header$ to $result$
            \State append $S[i]$ to $result$
            \State $i \gets i + run\_length$
        \Else \Comment{Literal run - collect up to 128 literals}
            \State $literal\_start \gets i$
            \State $literal\_count \gets 0$
            \While {$i \leq n$ and $literal\_count < 128$}
                \State \Comment{Check if next 3+ bytes form a run}
                \If {$i+2 \leq n$ and $S[i] = S[i+1]$ and $S[i] = S[i+2]$}
                    \State break \Comment{Found a run, stop collecting literals}
                \EndIf
                \State $i \gets i + 1$
                \State $literal\_count \gets literal\_count + 1$
            \EndWhile
            \State $header \gets literal\_count - 1$ \Comment{Positive header}
            \State append $header$ to $result$
            \For {$j \gets literal\_start$ to $literal\_start + literal\_count - 1$}
                \State append $S[j]$ to $result$
            \EndFor
        \EndIf
    \EndWhile
    \State \Return $result$
\EndProcedure
\end{algorithmic}
\end{algorithm}

\subsubsection{Decoding Algorithm}

\begin{algorithm}[H]
\caption{PackBits Decoding}
\begin{algorithmic}[1]
\Procedure{PackBits-Decode}{$data[1..m]$}
    \State $i \gets 1$
    \State $result \gets$ empty list
    \While {$i \leq m$}
        \State $h \gets data[i]$ \Comment{header byte as signed integer}
        \State $i \gets i + 1$
        \If {$h \ge 0$} \Comment{Literal run}
            \State $len \gets h + 1$
            \For {$j \gets 1$ to $len$}
                \State append $data[i]$ to $result$
                \State $i \gets i + 1$
            \EndFor
        \ElsIf {$h = -128$} \Comment{Reserved, usually skip}
            \State \Comment{No operation, sometimes used as marker}
        \Else \Comment{Repeated run ($h$ is negative)}
            \State $len \gets 1 - h$ \Comment{Since $h$ is negative}
            \State $value \gets data[i]$
            \State $i \gets i + 1$
            \For {$j \gets 1$ to $len$}
                \State append $value$ to $result$
            \EndFor
        \EndIf
    \EndWhile
    \State \Return $result$
\EndProcedure
\end{algorithmic}
\end{algorithm}

\subsubsection{Worked Examples}

\begin{examplebox}
\textbf{PackBits Example 1: Mixed Data}

Input: \texttt{AAAAABCDEFFFFGG}

\textbf{Encoding process:}
\begin{itemize}
    \item Run of 5 A's → header = -4 (since 1-(-4)=5), then byte 'A'
    \item Literals B,C,D,E (4 bytes) → header = 3 (copy 3+1=4 bytes), then B C D E
    \item Run of 4 F's → header = -3 (repeat 4 times), then byte 'F'
    \item Literals G,G (2 bytes) → header = 1 (copy 2 bytes), then G G
\end{itemize}

Encoded stream (as bytes with decimal headers):
\begin{verbatim}
[-4, 'A', 3, 'B','C','D','E', -3, 'F', 1, 'G','G']
\end{verbatim}

\textbf{Decoding:}
\begin{itemize}
    \item Read -4 → repeat next byte 5 times → output "AAAAA"
    \item Read 3 → copy next 4 bytes → output "BCDE"
    \item Read -3 → repeat next byte 4 times → output "FFFF"
    \item Read 1 → copy next 2 bytes → output "GG"
\end{itemize}

Result matches original!
\end{examplebox}

\begin{examplebox}
\textbf{PackBits Example 2: Long Literal Run}

Input: \texttt{ABCDEFGHIJKLMNOP} (16 distinct letters)

Since PackBits literal chunks max at 128 bytes, we can encode as one chunk:
\begin{itemize}
    \item Header = 15 (copy 16 bytes)
    \item Then all 16 letters
    \item Total: 1 header + 16 literals = 17 bytes
    \item Original: 16 bytes → slight expansion (1 byte)
\end{itemize}
\end{examplebox}

\begin{examplebox}
\textbf{PackBits Example 3: Long Repeated Run}

Input: 150 bytes of 'X'

We need multiple chunks:
\begin{itemize}
    \item First chunk: header = -127 (repeat 128 times), byte 'X'
    \item Second chunk: header = -21 (repeat 22 times), byte 'X'
    \item Total: 2 headers + 2 bytes = 4 bytes
    \item Original: 150 bytes
    \item Compression ratio: 37.5:1 (97.3\% savings!)
\end{itemize}
\end{examplebox}

\begin{examplebox}
\textbf{PackBits Example 4: Boundary Cases}

\textbf{Single byte:} Input: \texttt{X}
\begin{itemize}
    \item Must encode as literal
    \item Header = 0 (copy 1 byte), then 'X'
    \item 2 bytes for 1 byte → 100\% expansion
\end{itemize}

\textbf{Two identical bytes:} Input: \texttt{XX}
\begin{itemize}
    \item Option 1: Encode as run of 2 → header = -1 (repeat 2 times), then 'X'
    \item Option 2: Encode as literals → header = 1 (copy 2 bytes), then 'X','X'
    \item Both use 3 bytes vs original 2 → expansion
    \item PackBits has no efficient encoding for pairs!
\end{itemize}

\textbf{Two different bytes:} Input: \texttt{XY}
\begin{itemize}
    \item Literals: header = 1, then 'X','Y' → 3 bytes vs 2 → expansion
\end{itemize}
\end{examplebox}

\subsubsection{Optimal Use of PackBits}

\begin{importantbox}
\textbf{When PackBits Excels:}

\begin{itemize}
    \item \textbf{Long runs ($\geq$3):} Excellent compression
    \item \textbf{Mixed data with some long runs:} Good balance
    \item \textbf{Data with runs $\geq$128:} Still efficient (splits automatically)
\end{itemize}

\textbf{When PackBits Struggles:}

\begin{itemize}
    \item \textbf{Short runs (pairs):} No efficient encoding (3 bytes for 2)
    \item \textbf{Alternating data:} All literals, slight expansion
    \item \textbf{Very short chunks:} Header overhead dominates
\end{itemize}

\textbf{Break-even analysis:}
\begin{itemize}
    \item Literal chunk: cost = 1 + L bytes for L bytes (overhead = 1)
    \item Run chunk: cost = 2 bytes for L bytes (overhead = 2/L)
    \item For L=3, both cost 3 bytes
    \item For L=4, run saves 1 byte vs literals
    \item For L=2, run costs 3 bytes vs 2 bytes literals → worse
\end{itemize}
\end{importantbox}

\subsubsection{PackBits in Real-World Formats}

\begin{importantbox}
\textbf{Applications of PackBits:}

\begin{itemize}
    \item \textbf{TIFF:} Optional compression method (tag 32773)
    \item \textbf{Macintosh resource forks:} Used extensively
    \item \textbf{PDF:} Can use PackBits for stream compression
    \item \textbf{PostScript:} Supported as a filter
    \item \textbf{Many embedded systems:} Simple, efficient implementation
\end{itemize}

\textbf{Why PackBits remains popular:}
\begin{itemize}
    \item Simple to implement correctly
    \item No patents
    \item Works well on mixed data
    \item Byte-aligned (fast)
    \item Reasonable worst-case behavior (0.5-1 byte overhead per chunk)
\end{itemize}
\end{importantbox}

\newpage
\subsection{Zero Run-Length Encoding (Zero RLE)}

Zero RLE is a specialized variant that focuses exclusively on compressing runs of zero values. It's particularly important in transform-based compression systems.

\subsubsection{Motivation: Sparse Data in Transform Coding}

\begin{importantbox}
\textbf{Why Zero Runs Occur in Transform Coding:}

\begin{enumerate}
    \item \textbf{Energy compaction:} Transforms (DCT, wavelet) concentrate signal energy into few coefficients
    \item \textbf{Quantization:} Small coefficients become zero after quantization
    \item \textbf{High-frequency coefficients:} Often near-zero in natural images
    \item \textbf{Zigzag scanning:} Orders coefficients from low to high frequency, creating long zero runs at the end
\end{enumerate}
\end{importantbox}

\begin{definitionbox}
\textbf{Zero Run-Length Encoding:} A specialized form of RLE that focuses exclusively on runs of zero values, often combined with encoding the non-zero values that interrupt them.
\end{definitionbox}

\subsubsection{Basic Zero RLE Algorithm}

The simplest form of Zero RLE treats the data as a sequence of (zero run length, non-zero value) pairs.

\begin{algorithm}[H]
\caption{Zero RLE Encoding}
\begin{algorithmic}[1]
\Procedure{ZeroRLE-Encode}{$S[1..n]$}
    \State $result \gets$ empty list
    \State $i \gets 1$
    \While {$i \leq n$}
        \If {$S[i] = 0$}
            \State $zero\_run \gets 0$
            \While {$i \leq n$ and $S[i] = 0$}
                \State $zero\_run \gets zero\_run + 1$
                \State $i \gets i + 1$
            \EndWhile
            \If {$i > n$} \Comment{End-of-block special case}
                \State append $(zero\_run, \text{EOB})$ to $result$
            \Else
                \State append $(zero\_run, S[i])$ to $result$
                \State $i \gets i + 1$
            \EndIf
        \Else
            \State append $(0, S[i])$ to $result$ \Comment{Zero run length 0}
            \State $i \gets i + 1$
        \EndIf
    \EndWhile
    \State \Return $result$
\EndProcedure
\end{algorithmic}
\end{algorithm}

\subsubsection{Zero RLE in JPEG: A Detailed Case Study}

JPEG's encoding of AC coefficients is the most famous example of Zero RLE.

\begin{importantbox}
\textbf{JPEG's Zero-Run Encoding:}

After DCT and quantization, JPEG processes AC coefficients as follows:
\begin{enumerate}
    \item \textbf{Zigzag scan:} Reorder 8×8 block in zigzag pattern to maximize zero runs
    \item \textbf{Run-length encoding:} Encode as (RUNLENGTH, LEVEL) pairs
    \begin{itemize}
        \item \textbf{RUNLENGTH:} Number of consecutive zeros preceding this coefficient (0-15)
        \item \textbf{LEVEL:} Non-zero coefficient value
        \item Special case: (0,0) = End-of-Block (EOB)
        \item Special case: (15,0) = ZRL (16 zeros, handled specially)
    \end{itemize}
    \item \textbf{Huffman coding:} Each (RUNLENGTH, LEVEL) pair is Huffman coded
\end{enumerate}
\end{importantbox}

\begin{examplebox}
\textbf{JPEG Zigzag Scan Example:}

Original 8×8 block (quantized coefficients):
\begin{verbatim}
[150, -25, 0,   0,   0,   0,   12,  0
 0,   0,   0,   0,   0,   0,   0,   0
 0,   0,   0,   0,   0,   0,   0,   0
 0,   0,   0,   0,   0,   0,   0,   0
 0,   0,   0,   0,   0,   0,   0,   0
 0,   0,   0,   0,   0,   0,   0,   0
 0,   0,   0,   0,   0,   0,   0,   0
 0,   0,   0,   0,   0,   0,   0,   0]
\end{verbatim}

\textbf{Zigzag scan order (simplified):}
\begin{itemize}
    \item Position 1: 150 (DC coefficient - encoded separately)
    \item Positions 2-10: -25, 0, 0, 0, 0, 12, 0, 0, 0
    \item Remaining 54 positions: all zeros
\end{itemize}

\textbf{Zero-run encoding:}
\begin{itemize}
    \item After DC, first AC: (0, -25) meaning 0 zeros then -25
    \item Next: (4, 12) meaning 4 zeros then 12
    \item Then: (0,0) EOB meaning rest are zeros
\end{itemize}

\textbf{JPEG representation:} (0,-25), (4,12), EOB
\end{examplebox}

\begin{examplebox}
\textbf{JPEG Encoding Example:}

Consider this AC coefficient sequence:
0, 0, 0, 5, 0, 0, -3, 0, 0, 0, 0, 2, [rest zeros]

\textbf{Step 1: Identify zero runs and non-zero values:}
\begin{itemize}
    \item 3 zeros then 5
    \item 2 zeros then -3
    \item 4 zeros then 2
    \item All remaining zeros
\end{itemize}

\textbf{Step 2: Encode as (RUNLENGTH, LEVEL):}
\begin{itemize}
    \item (3,5)
    \item (2,-3)
    \item (4,2)
    \item (0,0) for EOB
\end{itemize}

\textbf{Step 3: Huffman coding:} Each pair is Huffman coded using tables optimized for typical distributions. The Huffman tables assign shorter codes to common (RUNLENGTH, LEVEL) combinations.
\end{examplebox}

\subsubsection{Handling Long Zero Runs}

JPEG's RUNLENGTH field is limited to 15, so special handling is needed for longer zero runs.

\begin{importantbox}
\textbf{JPEG's ZRL (Zero Run Length) Code:}

When a zero run exceeds 15 zeros:
\begin{itemize}
    \item Output (15,0) for each complete block of 16 zeros
    \item Then encode the remainder normally
\end{itemize}

\textbf{Example:} 37 zeros followed by a value 5
\begin{itemize}
    \item 37 zeros = 16 + 16 + 5
    \item Output: (15,0), (15,0), (5,5)
\end{itemize}

This allows arbitrary-length zero runs while keeping the RUNLENGTH field small.
\end{importantbox}

\subsubsection{Zero RLE in Other Transform Codecs}

\begin{importantbox}
\textbf{Zero RLE in Modern Codecs:}

\textbf{MPEG video codecs:}
\begin{itemize}
    \item Similar to JPEG for intra-coded blocks
    \item (run, level) coding for DCT coefficients
    \item Variable-length codes optimized for video statistics
\end{itemize}

\textbf{JPEG 2000 (wavelet-based):}
\begin{itemize}
    \item Embedded block coding with optimized truncation (EBCOT)
    \item Fractional bit-plane coding
    \item Zero-run encoding implicit in bit-plane representation
    \item Arithmetic coding with context modeling
\end{itemize}

\textbf{H.264/AVC:}
\begin{itemize}
    \item CAVLC (Context-Adaptive Variable-Length Coding)
    \item Encodes runs of zeros before non-zero coefficients
    \item Adapts to local statistics
\end{itemize}
\end{importantbox}

\subsubsection{Zero RLE Pipeline in Transform Coding}

\begin{importantbox}
\textbf{Complete Zero-RLE Pipeline:}

\begin{verbatim}
Raw Image → Transform → Quantization → Zigzag Scan → 
Zero-RLE → Entropy Coding → Compressed Bitstream
\end{verbatim}

\textbf{Why this pipeline works:}
\begin{enumerate}
    \item \textbf{Transform:} Converts spatial redundancy to frequency coefficients
    \item \textbf{Quantization:} Reduces precision, creates zeros
    \item \textbf{Zigzag scan:} Reorders coefficients to maximize zero runs
    \item \textbf{Zero-RLE:} Compresses zero runs efficiently
    \item \textbf{Entropy coding:} Further compresses (run, level) pairs
\end{enumerate}
\end{importantbox}

\newpage
\subsection{Comparison of RLE Variants}

\subsubsection{Decision Matrix}

\begin{tabular}{|p{2.5cm}|p{2.6cm}|p{2.6cm}|p{2.6cm}|p{2.6cm}|p{2.6cm}|}
\hline
\textbf{Criteria} & \textbf{Count-Based Fixed} & \textbf{Count-Based Variable} & \textbf{Escape-Based} & \textbf{PackBits} & \textbf{Zero RLE} \\
\hline
Implementation complexity & Very low & High & Medium & Medium & Medium \\
\hline
Speed & Very fast & Slow & Fast & Fast & Fast \\
\hline
Random access & Yes & No & No & No & No \\
\hline
Best for & Uniform runs & Skewed runs & Mixed data & Mixed data & Sparse data \\
\hline
Worst-case expansion & 100\% & ~100\% (plus bit overhead) & Small (escape handling) & Small (header overhead) & Small \\
\hline
Max run length & Fixed limit & Unlimited & Unlimited & 128 per chunk & Unlimited \\
\hline
Common uses & Simple formats, hardware & Advanced codecs, BWT pipeline & TIFF, some fax & TIFF, Mac formats & JPEG, MPEG, DCT-based codecs \\
\hline
\end{tabular}

\subsubsection{When to Use Each Variant}

\begin{importantbox}
\textbf{Selection Guide:}

\begin{itemize}
    \item \textbf{Choose Fixed-Width Count RLE when:}
    \begin{itemize}
        \item You need maximum speed and simplicity
        \item Run lengths are bounded and relatively uniform
        \item Hardware implementation is required
        \item Random access to runs is needed
    \end{itemize}
    
    \item \textbf{Choose Variable-Width Count RLE when:}
    \begin{itemize}
        \item Run length distribution is highly skewed (many short runs)
        \item Maximum compression ratio is critical
        \item You can tolerate bit-level I/O complexity
        \item Part of a larger pipeline (e.g., BWT+MTF)
    \end{itemize}
    
    \item \textbf{Choose Escape-Based RLE when:}
    \begin{itemize}
        \item Data has mixed runs and literals
        \item You need to guarantee no worst-case expansion
        \item Simple byte-aligned implementation is desired
    \end{itemize}
    
    \item \textbf{Choose PackBits when:}
    \begin{itemize}
        \item You need a standard, widely-supported format
        \item Data has moderate runs and literals
        \item Implementation simplicity is important
        \item Maximum run length of 128 is acceptable
    \end{itemize}
    
    \item \textbf{Choose Zero RLE when:}
    \begin{itemize}
        \item Working with transform-coded data (images, video)
        \item Data is sparse with many zeros
        \item Part of a larger compression pipeline
        \item Non-zero values need special handling
    \end{itemize}
\end{itemize}
\end{importantbox}

\subsection{Limitations of Run-Length Encoding}

\subsubsection{Why RLE Fails on Natural Language Text}

Natural language text represents RLE's worst nightmare.

\begin{importantbox}
\textbf{Characteristics of Natural Language Text:}

\begin{enumerate}
    \item \textbf{Large alphabet:} 26 letters + punctuation + case + spaces
    \item \textbf{Rare repeats:} In English, repeated letters are uncommon
    \begin{itemize}
        \item Double letters: "book", "little", "success" — relatively rare
        \item Triple letters: Almost nonexistent (except "beeeeep" in comics)
    \end{itemize}
    \item \textbf{Statistical distribution:} Letter frequencies follow Zipf's law
    \item \textbf{Context structure:} Redundancy is in patterns, not runs
\end{enumerate}
\end{importantbox}

\begin{examplebox}
\textbf{RLE on Shakespeare's Sonnet:}

Consider the opening of Sonnet 18:
\texttt{Shall I compare thee to a summer's day?}

\textbf{Run analysis:}
\begin{itemize}
    \item Total characters (including spaces): ~500
    \item Runs of length 1: ~490
    \item Runs of length 2: ~5 (double letters: "ll" in "shall", "mm" in "summer")
    \item Runs of length 3: 0
    \item Average run length: ~1.01
\end{itemize}

\textbf{RLE encoding:}
\begin{itemize}
    \item Original: 500 bytes
    \item Encoded: ~500 runs × 2 = 1000 values
    \item At 1 byte per value: 1000 bytes (100\% expansion)
\end{itemize}

\textbf{Why the expansion?}
\begin{itemize}
    \item Each run of length 1 requires a count of 1
    \item The count field adds 100\% overhead to each character
    \item The few runs of length 2 save only 1 byte each (2 chars → 2 bytes vs 4 bytes for two singles)
    \item Savings from runs: 5 runs × (2 bytes saved per run) = 10 bytes saved
    \item But overhead from counts: 500 extra bytes
    \item Net loss: 490 bytes
\end{itemize}
\end{examplebox}

\subsubsection{Sensitivity to Symbol Ordering}

RLE's effectiveness depends critically on how symbols are arranged.

\begin{definitionbox}
\textbf{Order Sensitivity:} The property that RLE compression ratio can change dramatically if the same data is reordered, even though the multiset of symbols remains identical.
\end{definitionbox}

\begin{examplebox}
\textbf{Dramatic Effect of Ordering:}

Consider the multiset: 100 A's, 100 B's, 100 C's

\textbf{Ordering 1: Interleaved (ABC ABC ABC...)} 
\texttt{ABCABCABC...} (300 bytes)
\begin{itemize}
    \item Each run length = 1
    \item Number of runs = 300
    \item Encoded size: 600 values (assuming 1-byte counts)
    \item Result: 100\% expansion!
\end{itemize}

\textbf{Ordering 2: Grouped (AAA...BBB...CCC...)}
\texttt{AAA... (100) BBB... (100) CCC... (100)}
\begin{itemize}
    \item 3 runs, each length 100
    \item Encoded size: 3×2 = 6 values
    \item Result: 98\% compression!
\end{itemize}
\end{examplebox}

\subsection{RLE as a Preprocessing Technique}

\subsubsection{Combining RLE with Huffman Coding}

The RLE+Huffman combination is a classic and effective approach.

\begin{importantbox}
\textbf{Why Combine RLE and Huffman:}

\begin{itemize}
    \item \textbf{RLE handles runs:} Converts repeated symbols into (symbol, count) pairs
    \item \textbf{Huffman handles frequencies:} Assigns shorter codes to frequent symbol/count combinations
    \item \textbf{Synergy:} Each addresses a different type of redundancy
\end{itemize}
\end{importantbox}

\begin{examplebox}
\textbf{CCITT Group 3 Fax: The Classic Example}

Group 3 fax does exactly this:
\begin{itemize}
    \item Step 1: Identify runs of white and black
    \item Step 2: For each run length, use a Huffman code
    \item Step 3: Special make-up codes for long runs
    \item Step 4: Separate Huffman tables for white and black runs
\end{itemize}

\textbf{Result:} Achieves ~10-20:1 compression on typical documents.
\end{examplebox}

\subsubsection{The Crucial Role of RLE in BWT Pipeline}

This is the key connection to the next chapter.

\begin{importantbox}
\textbf{The BWT Pipeline:}

\begin{verbatim}
Original Text → BWT → MTF → RLE → Huffman/Arithmetic
\end{verbatim}

\textbf{Why RLE is essential after BWT+MTF:}

\begin{enumerate}
    \item \textbf{BWT groups similar contexts:} Characters with similar contexts end up near each other
    \item \textbf{MTF converts to small numbers:} Recently seen characters get small indices (often 0)
    \item \textbf{Result:} Long runs of zeros and other small numbers appear
    \item \textbf{RLE compresses these runs:} Exactly what RLE is good at!
\end{enumerate}
\end{importantbox}

\subsection{Summary}

\begin{importantbox}
\textbf{Key Takeaways:}

\begin{enumerate}
    \item \textbf{Equality redundancy} is the simplest form of data redundancy—consecutive identical symbols.
    
    \item \textbf{Four major RLE variants} address different data characteristics:
    \begin{itemize}
        \item \textit{Count-Based RLE:} Classic (symbol, count) approach. Fixed-width for simplicity, variable-width for optimal compression.
        \item \textit{Escape-Based RLE:} Uses markers to distinguish runs from literals. Prevents worst-case expansion.
        \item \textit{Block-Based RLE (PackBits):} Header-based chunks balance overhead and efficiency.
        \item \textit{Zero RLE:} Specialized for sparse data, essential in transform coding.
    \end{itemize}
    
    \item \textbf{Performance bounds:}
    \begin{itemize}
        \item Best case: near-infinite compression (single run)
        \item Worst case: 100\% expansion (no runs, fixed-width)
        \item Break-even: depends on count field size and run length distribution
    \end{itemize}
    
    \item \textbf{RLE is passive:} It can only exploit existing runs. This limitation motivates transform-based compression (BWT) which \textit{creates} runs.
\end{enumerate}
\end{importantbox}

\begin{tcolorbox}[colback=orange!10!white,colframe=orange!50!black, title=Preview: The Burrows-Wheeler Transform]
\begin{center}
\textbf{From RLE to BWT: The Evolution of Compression Thinking}

\vspace{0.3cm}
\textbf{Level 1:} "I see runs, I compress runs." (Basic RLE)

\vspace{0.2cm}
\textbf{Level 2:} "I adapt RLE to handle different situations." (RLE variants)

\vspace{0.2cm}
\textbf{Level 3:} "I transform data to create the runs RLE needs." (BWT)

\vspace{0.3cm}
The next chapter will show you how to achieve Level 3 thinking.
\end{center}
\end{tcolorbox}

\end{document}
